\begin{figure}
\begin{minipage}[t]{0.08\textwidth}
\vspace{25pt}
\textbf{$f(x)$} \\
\\[31pt]
\textbf{$g_{1.25}(x)$} \\
\\[31pt]
\textbf{$g_{1.5}(x)$} \\
\\[31pt]
\textbf{$f(x)$} \\
\\[31pt]
\textbf{$g_{1.25}(x)$} \\
\\[31pt]
\textbf{$g_{1.5}(x)$}
\end{minipage}%
\begin{minipage}[t]{0.92\textwidth}
\vspace{0pt}
\centering  
\addtolength{\tabcolsep}{-5.0pt}
\begin{tabular}{c}
\begin{overpic}[width=\ResultFigWidth]{\GridFig{Pyramid}{9}{mala}}\end{overpic} \\
\begin{overpic}[width=\ResultFigWidth]{\GridErrFig{Pyramid}{9}{mala}}
    \put(14.3, 0){\begin{tikzpicture}[x=1pt,y=1pt]
        \draw[red,ultra thick] (0,0) rectangle (61.5,190.2);
    \end{tikzpicture}}
    \put(28.7, 0){\begin{tikzpicture}[x=1pt,y=1pt]
        \draw[red,ultra thick] (0,0) rectangle (61.5,190.2);
    \end{tikzpicture}}
    \put(43, 0){\begin{tikzpicture}[x=1pt,y=1pt]
        \draw[red,ultra thick] (0,0) rectangle (62,190.2);
    \end{tikzpicture}}
\end{overpic} \\
\end{tabular}
\end{minipage}
\caption{Grid9Pyramid-MALA 前三行为render的结果，后三行为error map。最后一列为使用的reference原图。通过error map的第二三四列都可以看出，$f(x)$的error map的区域间相对亮度误差更大。}
\label{fig:Grid9PyramidFigMALA}
\end{figure}

\begin{figure}
\begin{minipage}[t]{0.08\textwidth}
\vspace{25pt}
\textbf{$f(x)$} \\
\\[31pt]
\textbf{$g_{1.25}(x)$} \\
\\[31pt]
\textbf{$g_{1.5}(x)$} \\
\\[31pt]
\textbf{$f(x)$} \\
\\[31pt]
\textbf{$g_{1.25}(x)$} \\
\\[31pt]
\textbf{$g_{1.5}(x)$}
\end{minipage}%
\begin{minipage}[t]{0.92\textwidth}
\vspace{0pt}
\centering  
\addtolength{\tabcolsep}{-5.0pt}
\begin{tabular}{c}
\begin{overpic}[width=\ResultFigWidth]{\GridFig{Sphere}{100}{h2mc}}\end{overpic} \\
\begin{overpic}[width=\ResultFigWidth]{\GridErrFig{Sphere}{100}{h2mc}}
    \put(7.4, 34.2){\begin{tikzpicture}[x=1pt,y=1pt]
        \draw[red, thick] (0,0) rectangle (10,10);
    \end{tikzpicture}}
    \put(21.8, 34.2){\begin{tikzpicture}[x=1pt,y=1pt]
        \draw[red, thick] (0,0) rectangle (10,10);
    \end{tikzpicture}}
    \put(36.2, 34.2){\begin{tikzpicture}[x=1pt,y=1pt]
        \draw[red, thick] (0,0) rectangle (10,10);
    \end{tikzpicture}}
    \put(38.7, 31.7){\begin{tikzpicture}[x=1pt,y=1pt]
        \draw[red, thick] (0,0) rectangle (10,10);
    \end{tikzpicture}}
    \put(50.6, 34.2){\begin{tikzpicture}[x=1pt,y=1pt]
        \draw[red, thick] (0,0) rectangle (10,10);
    \end{tikzpicture}}
    \put(53.1, 31.7){\begin{tikzpicture}[x=1pt,y=1pt]
        \draw[red, thick] (0,0) rectangle (10,10);
    \end{tikzpicture}}
    \put(65, 34.2){\begin{tikzpicture}[x=1pt,y=1pt]
        \draw[red, thick] (0,0) rectangle (10,10);
    \end{tikzpicture}}
    \put(11.3, 16.2){\begin{tikzpicture}[x=1pt,y=1pt]
        \draw[red, thick] (0,0) rectangle (10,10);
    \end{tikzpicture}}
    \put(25.7, 16.2){\begin{tikzpicture}[x=1pt,y=1pt]
        \draw[red, thick] (0,0) rectangle (10,10);
    \end{tikzpicture}}
\end{overpic} \\
\end{tabular}
\end{minipage}
\caption{Grid100Sphere-H2MC 前三行为render的结果，后三行为error map。最后一列为使用的reference原图。通过error map可以看出，$g_{1.25}(x)$和$g_{1.5}(x)$相较于$f(x)$，各个区域亮度分配更均匀。}
\label{fig:Grid100SphereFigH2MC} 
\end{figure}

\begin{figure}
\begin{minipage}[t]{0.08\textwidth}
\vspace{25pt}
\textbf{$f(x)$} \\
\\[31pt]
\textbf{$g_{1.25}(x)$} \\
\\[31pt]
\textbf{$g_{1.5}(x)$} \\
\\[31pt]
\textbf{$f(x)$} \\
\\[31pt]
\textbf{$g_{1.25}(x)$} \\
\\[31pt]
\textbf{$g_{1.5}(x)$}
\end{minipage}%
\begin{minipage}[t]{0.92\textwidth}
\vspace{0pt}
\centering  
\addtolength{\tabcolsep}{-5.0pt}
\begin{tabular}{c}
\begin{overpic}[width=\ResultFigWidth]{\GridFig{Teapot}{25}{mala}}\end{overpic} \\
\begin{overpic}[width=\ResultFigWidth]{\GridErrFig{Teapot}{25}{mala}}\end{overpic} \\
\end{tabular}
\end{minipage}
\caption{Grid25Teapot-MALA 前三行为render的结果，后三行为error map。最后一列为使用的reference原图。这一场景中茶壶来自veachdoor中的mesh，从MSE来看，$g(x)$仍然优于$f(x)$的结果。}
\label{fig:Grid25TeapotFigMALA}
\end{figure}
